% preface_sections2_4_draft.tex — Sections 2-4 of the Restored Preface
% To be integrated into preface.tex
% Macros assumed from main.tex: \cA, \barB, \gAmod, \MC, \cZ
\providecommand{\cA}{\mathcal{A}}
\providecommand{\barB}{\overline{B}}
\providecommand{\gAmod}{\mathfrak{g}_{\cA}^{\mathrm{mod}}}
\providecommand{\MC}{\text{MC}}
\providecommand{\cM}{\mathcal{M}}
\providecommand{\cC}{\mathcal{C}}
\providecommand{\cO}{\mathcal{O}}
\providecommand{\cP}{\mathcal{P}}
\providecommand{\Eone}{E_1}
\providecommand{\Convinf}{\mathrm{Conv}_\infty}
\providecommand{\Convstr}{\mathrm{Conv}_{\mathrm{str}}}
\providecommand{\orline}[1]{\mathrm{or}(#1)}
\providecommand{\fg}{\mathfrak{g}}

% ====================================================================
\section*{2.\quad Curvature and the genus tower}
% ====================================================================

\medskip

\noindent
On a curve of genus $g \geq 1$, the bar differential does not square
to zero.  The propagator $d\log(z_1-z_2)$ has no global analogue:
analytic continuation around $B$-cycles introduces period corrections,
these corrections break the Arnold relation, and the failure is
measured by a single scalar $\kappa(\cA)$ extracted from the leading
OPE singularity.  The problem is to restore nilpotence at all genera
and to determine the genus-$g$ obstruction class.  The answer is that
period integrals absorb the curvature, and the obstruction class is
$\kappa(\cA) \cdot \lambda_g$ with $\hat A$-genus generating function.

\subsection*{The genus-$g$ propagator}

At genus $g\ge 1$, the difference $z_1-z_2$ has no global meaning.
The replacement is the \emph{prime form} $E(z_1,z_2)$: a section
of $K^{-1/2}\boxtimes K^{-1/2}$ on $X\times X$, vanishing to
first order along the diagonal $\Delta$, with no other zeros.  The
bundle is $K^{-1/2}\boxtimes K^{-1/2}$, \emph{not}
$K^{+1/2}\boxtimes K^{+1/2}$: the prime form is a
$(-\frac12,-\frac12)$-differential.  In a local coordinate where
the diagonal is $z_1=z_2$,
\[
E(z_1,z_2)
\;=\;
(z_1-z_2)\cdot\bigl(1+O(z_1-z_2)\bigr)
\cdot(dz_1)^{-1/2}(dz_2)^{-1/2}.
\]

The prime form has nontrivial monodromy.  Analytic continuation
around the $B$-cycle $B_j$ in the variable~$z_1$:
\[
E(z_1+B_j,z_2)
\;=\;
-\exp\bigl(
-\pi i\Omega_{jj}-2\pi i\textstyle\int_{z_2}^{z_1}\omega_j
\bigr)
\cdot E(z_1,z_2),
\]
where $\Omega=(\Omega_{jk})$ is the period matrix of~$X$
($\Omega_{jk}=\oint_{B_k}\omega_j$) and
$\omega_1,\dots,\omega_g$ are the holomorphic differentials
normalised by $\oint_{A_j}\omega_k=\delta_{jk}$.  The
$A$-cycle monodromy is trivial: $E(z_1+A_j,z_2)=E(z_1,z_2)$.

The logarithmic derivative
$\omega(z_1,z_2) := d_{z_1}\log E(z_1,z_2)$
is the fundamental meromorphic differential of the second kind:
symmetric, a double pole on~$\Delta$, vanishing $A$-periods.  The
\emph{genus-$g$ propagator}:
\[
\eta^{(g)}_{12}
\;=\;
d\log E(z_1,z_2)
\;+\;
\pi\sum_{j,k=1}^{g}
(\mathrm{Im}\,\Omega)^{-1}_{jk}\,
\mathrm{Im}\!\Bigl(\int_{z_2}^{z_1}\omega_j\Bigr)\,
\omega_k(z_1).
\]
The period correction $(\mathrm{Im}\,\Omega)^{-1}_{jk}$
compensates the $B$-cycle monodromy, rendering $\eta^{(g)}_{12}$
single-valued on~$X$.  The price: $\eta^{(g)}_{12}$ is no longer
holomorphic; the correction term is a $(1,0)$-form in~$z_1$
multiplied by a real-analytic function of both variables.  The
propagator couples to the Hodge bundle
$\mathbb E\to\overline{\cM}_g$ through the period
matrix~$\Omega$.

The propagator $d\log E(z_1,z_2)$ has conformal weight~$1$ in both
variables, regardless of the conformal weight~$h$ of the field
being sewed.  Every edge of every genus-$g$ graph sum uses the
standard Hodge bundle $\mathbb E_1 = R^0\pi_*\omega_\pi$.  The
Mumford isomorphism $c_1(\mathbb E_j) = (6j^2-6j+1)\lambda_1$
would produce discrepancies of $13\times$ (Virasoro) or $22.6\times$
(W$_3$) at genus~$1$ if weight-$h$ generators were assigned to
$\mathbb E_h$: the propagator argument eliminates this.

\subsection*{Curvature}

The Arnold relation at genus~$g$ acquires a defect.  At genus~$0$,
\[
\eta_{12}\wedge\eta_{23}
+\eta_{23}\wedge\eta_{31}
+\eta_{31}\wedge\eta_{12}
\;=\;0
\]
holds identically.  At genus $g\ge 1$, the period corrections
produce a residual $(1,1)$-form: the three-term sum does not vanish.
The fibrewise bar differential satisfies
\begin{equation}\label{eq:pref-curvature}
d_{\mathrm{fib}}^{\,2}\;=\;\kappa(\cA)\cdot\omega_g,
\end{equation}
where $\omega_g$ is the Arakelov $(1,1)$-form on the fibre of the
universal curve $\pi\colon\mathcal C_g\to\overline{\cM}_g$
and $\kappa(\cA)\in\Bbbk$ is the scalar \emph{modular
characteristic}, determined entirely by the leading OPE singularity.

For the Heisenberg algebra: the only OPE datum entering the bar
differential is $\alpha_{(1)}\alpha=k$.  Every residue at a
codimension-one collision stratum extracts this coefficient:
\[
d_{\mathrm{fib}}^{\,2}\big|_{\mathcal H_k}
\;=\;
k\cdot\omega_g,
\qquad\text{so}\quad
\kappa(\mathcal H_k)=k.
\]
For the Kac--Moody algebra: the bracket
$J^a_{(0)}J^b=f^{ab}{}_{c}J^c$ and the inner product
$J^a_{(1)}J^b=k\kappa^{ab}$ both contribute.  The curvature
receives the trace over the adjoint representation:
\[
\kappa(\widehat{\fg}_k)
\;=\;
\frac{(k+h^\vee)\dim\fg}{2h^\vee}\,.
\]
The numerator $k+h^\vee$ is the \emph{shifted level}; $h^\vee$ is
the dual Coxeter number.  At $k=-h^\vee$ (the critical level),
$\kappa=0$ and the bar complex is \emph{uncurved}: the Sugawara
construction is undefined (the central charge
$c=k\dim\fg/(k+h^\vee)$ has a pole), but the curvature
parameter vanishes.

For the Virasoro algebra at central charge~$c$: the stress tensor
$T$ has OPE $T(z)T(w)\sim\frac{c/2}{(z-w)^4}
+\frac{2T(w)}{(z-w)^2}+\frac{\partial T(w)}{z-w}$.
The bar propagator $d\log E(z,w)$ absorbs one pole order
(the $d\log$ measure sends $z^{-n}$ to $z^{-(n-1)}$), so the
collision residue $r(z)=(c/2)/z^3+2T/z$ has pole orders one
less than the OPE.  The curvature:
\[
\kappa(\mathrm{Vir}_c)\;=\;\frac{c}{2}\,.
\]
The Koszul dual is $\mathrm{Vir}_{26-c}$, with
$\kappa(\mathrm{Vir}_{26-c})=(26-c)/2$: the two curvatures
sum to~$13$, not zero.  At $c=26$ the dual is uncurved; at
$c=0$ the algebra itself is uncurved; at $c=13$ the algebra
is self-dual under Koszul duality.

\subsection*{Period integrals restore nilpotence}

The curvature~\eqref{eq:pref-curvature} obstructs
$d_{\mathrm{fib}}^2=0$, but not the total $D_g^2=0$.
Integration against $A$-cycle periods defines a connection on the
family of bar complexes parametrised by~$\overline{\cM}_g$.
The corrected total differential
\[
D_g\;=\;d_{\mathrm{fib}}+\nabla^{\mathrm{GM}}
\]
incorporates the Gauss--Manin connection on the Hodge bundle
(the canonical flat connection on $R^1\pi_*\mathbb C$,
encoding how the periods of the curve vary with moduli), and
$D_g^{\,2}=0$: the $(1,1)$-curvature of the fibrewise differential
is absorbed by the $(0,1)$-part of the base connection.

The mechanism: the fibrewise curvature $\kappa\cdot\omega_g$ is a
class in $H^{1,1}(\mathcal C_g)$; the period map
$\overline{\cM}_g\to\mathcal A_g$ identifies this class with the
$(0,1)$-curvature of the Gauss--Manin connection on
$R^1\pi_*\mathbb C$.  The two curvatures cancel:
\[
D_g^{\,2}
\;=\;
d_{\mathrm{fib}}^2
+[d_{\mathrm{fib}},\nabla^{\mathrm{GM}}]
+(\nabla^{\mathrm{GM}})^2
\;=\;
\kappa\cdot\omega_g-\kappa\cdot\omega_g+0
\;=\;0.
\]

\subsection*{The genus tower}

With $D_g^2 = 0$ restored, the bar complex defines a flat family of
cochain complexes over $\overline{\cM}_g$; the obstruction to
triviality of this family is a characteristic class.  For
uniform-weight algebras (single-generator, or multi-generator with
all generators of the same conformal weight), it takes the form
\[
\mathrm{obs}_g(\cA)
\;=\;
\kappa(\cA)\cdot\lambda_g
\;\in\;
H^*(\overline{\cM}_g),
\qquad
\lambda_g\;=\;c_g(\mathbb E),
\]
where $\mathbb E = \pi_*\omega_\pi$ is the \emph{Hodge bundle}:
the vector bundle on $\overline{\cM}_g$ whose fibre over a
point $[C]$ is the $g$-dimensional space
$H^0(C, \omega_C)$ of holomorphic differentials on~$C$.
At genus~$1$ this factorization holds unconditionally for all
families; for multi-weight families at genus~$g \geq 2$ the scalar
formula \emph{fails}: a cross-channel correction from
mixed-propagator boundary graphs appears
(Theorem~\ref{thm:multi-weight-genus-expansion}).

Grothendieck--Riemann--Roch applied to the universal curve
$\pi\colon\mathcal C_g\to\overline{\cM}_g$ computes the
total obstruction: the Chern character of $R\pi_*\omega_\pi$
involves the Todd class of the relative tangent bundle~$T_\pi$,
which is the $\hat A$-genus.  The generating function:
\[
\sum_{g\ge 1}\mathrm{obs}_g(\cA)\cdot x^{2g}
\;=\;
\kappa(\cA)\cdot\bigl(\hat A(ix)-1\bigr).
\]
Since $\hat A(ix)=(x/2)/\sin(x/2)
=1+x^2/24+7x^4/5760+\cdots$ (all positive),
the first Faber--Pandharipande coefficients are:
\begin{align*}
F_1(\cA)
&\;=\;
\frac{\kappa(\cA)}{24}\,,\\[3pt]
F_2(\cA)
&\;=\;
\frac{7\,\kappa(\cA)}{5760}\,,\\[3pt]
F_3(\cA)
&\;=\;
\frac{31\,\kappa(\cA)}{967680}\,.
\end{align*}
The genus-$g$ coefficient is
$\kappa(\cA)\cdot\tfrac{2^{2g-1}-1}{2^{2g-1}}
\cdot\tfrac{|B_{2g}|}{(2g)!}$.
The $\hat A$-genus appears from the bar complex as the connection
between chiral algebra and index theory.

The scalar $\kappa(\cA)$ is only the leading term.  The bar complex
carries higher-arity data: a universal Maurer--Cartan element
$\Theta_\cA$ whose projection to arity~$r$ yields a shadow invariant
$S_r(\cA)$.  On each primary deformation line~$L$, the shadow
generating function $H(t)=\sum_{r\ge 2}rS_r\,t^r$ satisfies
\begin{equation}\label{eq:preface-riccati-hook}
H(t)^2 \;=\; t^4\,Q_L(t),
\qquad
Q_L(t) = 4\kappa^2+12\kappa\alpha\,t+(9\alpha^2+16\kappa S_4)\,t^2,
\end{equation}
a quadratic polynomial in three genus-$0$ invariants
$(\kappa,\alpha,S_4)$
(Theorem~\ref{thm:riccati-algebraicity}).  The infinite tower is
algebraic of degree~$2$: determined, computable, and governed by the
single discriminant $\Delta = 8\kappa S_4$.

\subsection*{Complementarity}

The genus-$g$ cohomology of the center local system decomposes into
two Lagrangian halves: one controlled by~$\cA$, the other
by the Koszul dual~$\cA^!$.  Verdier duality on the Ran space
produces the splitting; shifted-symplectic geometry governs the
pairing:
\[
R\Gamma(\overline{\cM}_g,\;\mathcal Z_\cA)
\;\simeq\;
\mathbf Q_g(\cA)\;\oplus\;\mathbf Q_g(\cA^!).
\]
What~$\cA$ sees as deformation, $\cA^!$ sees as obstruction.
The natural pairing
\[
\mathbf Q_g(\cA)\otimes\mathbf Q_g(\cA^!)
\;\longrightarrow\;
R\Gamma(\overline{\cM}_g,\;\omega_{\overline{\cM}_g})
\;\simeq\;\Bbbk[-(3g-3)]
\]
is a $(-(3g-3))$-shifted symplectic form, and the two summands are
\emph{Lagrangian}: each is isotropic and of half the total
dimension.  At the scalar level, complementarity reduces to
\[
\kappa(\cA)+\kappa(\cA^!)\;=\;0
\qquad\text{(for Kac--Moody and free-field algebras).}
\]
For $\mathcal{W}$-algebras the sum $\kappa+\kappa^!$ is a nonzero
constant: $13$ for Virasoro, $250/3$ for~$\mathcal{W}_3$.

\subsection*{Feigin--Frenkel duality}

The Sugawara construction produces the Virasoro field
\[
T(z)
\;=\;
\frac{1}{2(k+h^\vee)}\sum_{a=1}^{\dim\fg}
{:}J^a(z)J^a(z){:}\,,
\]
with central charge $c=k\dim\fg/(k+h^\vee)$.  At the
critical level $k=-h^\vee$, the denominator vanishes: $T(z)$ is
\emph{undefined}, not ``$c$ diverges.''  The chiral algebra
acquires a large center, the Feigin--Frenkel center
$\mathfrak z(\widehat{\fg}_{-h^\vee})$, isomorphic to the
algebra of functions on the space of local ${}^L\fg$-opers
on the formal disc.

The Feigin--Frenkel involution
\[
k\;\longleftrightarrow\;-k-2h^\vee
\]
is Koszul duality on configuration spaces: it interchanges~$\cA$
and~$\cA^!$ while preserving the bar-cobar adjunction.  At the
level of modular characteristics,
$\kappa(\widehat{\fg}_k)+
\kappa(\widehat{\fg}_{-k-2h^\vee})=0$:
complementarity is a direct consequence.  At the critical
level, the Koszul dual is the Langlands-dual critical-level
algebra~$\widehat{{}^L\fg}_{-{}^Lh^\vee}$, not a
renormalised copy of the original.

\subsection*{The total space}

The universal configuration fibration
\[
\pi_n\colon
\overline{\mathcal C}_n(\mathcal C_g)
\;\longrightarrow\;
\overline{\cM}_g
\]
packages the bar complex into a single global object.  Its boundary
divisors carry two distinct geometric roles:
\begin{itemize}
\item \emph{Collision divisors} (vertical): the fibrewise boundary
of $\overline{\mathrm{FM}}_n(X_s)$ over a point
$s\in\overline{\cM}_g$.  These produce the bar
differential~$d_{\mathrm{res}}$, the OPE residues.
\item \emph{Degeneration divisors} (horizontal): the boundary of
$\overline{\cM}_g$ itself (nodal curves, reducible
fibres).  These produce the period corrections, the
Gauss--Manin connection terms in~$D_g$.
\end{itemize}
The collision divisors are the algebraic boundary (controlled by
the Borcherds identity); the degeneration divisors are the
geometric boundary (controlled by the period map and the
clutching morphisms of stable curves).  The interaction between them
is the curvature~$\kappa$: the fibrewise bar differential fails to
be nilpotent precisely because the vertical and horizontal
boundaries do not commute.


% ====================================================================
\section*{3.\quad Modular homotopy theory}
% ====================================================================

\medskip

\noindent
The fibrewise differential~$d_{\mathrm{fib}}$, the Gauss--Manin
connection~$\nabla^{\mathrm{GM}}$, and the clutching morphisms at
boundary strata of $\overline{\cM}_{g,n}$ satisfy compatibility
relations that no tree-level algebraic structure can encode: the
self-gluing of a curve at two marked points raises genus by one,
an operation with no tree-level counterpart.  The correct algebraic
structure is the \emph{modular operad} (Getzler--Kapranov):
operations indexed by connected stable graphs of arbitrary genus,
with both separating and non-separating gluing.  The bar complex of a
chiral algebra is an algebra over the Feynman transform of the
commutative modular operad; the total differential $D^2 = 0$ is
the codimension-two boundary cancellation on the moduli space of
stable curves, lifted through the log-FM residue correspondence.

\subsection*{3.1.\enspace The modular operad}

A \emph{stable graph}~$\Gamma$ parametrises a boundary stratum of
the Deligne--Mumford moduli space $\overline{\cM}_{g,n}$: vertices
are irreducible components, edges are nodes, genus labels record the
genera of the components, half-edges are marked points.

A \emph{modular operad}~$\cO$ has operations indexed
by \emph{connected stable graphs}.  A connected
stable graph~$\Gamma$ has a genus label
$g(v)\in\mathbb Z_{\ge 0}$ at each vertex satisfying
$2g(v)-2+\mathrm{val}(v)>0$ (stability), and total genus
$g(\Gamma)=b_1(\Gamma)+\sum_v g(v)$; trees are the
genus-$0$ case.
The structure maps of~$\cO$ are of two kinds:
\emph{separating gluing}
$\xi_{\mathrm{sep}} \colon
\cO(g_1, n_1{+}1) \otimes \cO(g_2, n_2{+}1)
\to \cO(g_1{+}g_2, n_1{+}n_2)$
(joining two components at a node), and
\emph{non-separating gluing}
$\xi_{\mathrm{ns}} \colon \cO(g, n{+}2)
\to \cO(g{+}1, n)$
(identifying two marked points on the same component,
raising genus by~$1$).  These satisfy $\Sigma_n$-equivariance,
associativity, and commutativity.
The set
$\mathsf{Gr}^{\mathrm{st}}_{g,n}$ parametrises the boundary
stratification of $\overline{\cM}_{g,n}$.  Every codimension-one
boundary divisor is a one-edge degeneration; every codimension-two
stratum is the intersection of exactly two codimension-one divisors,
appearing with opposite orientations in the boundary-of-boundary
complex.

\subsection*{3.2.\enspace The Feynman transform}

The \emph{Feynman transform} of~$\cO$, the modular analogue of
the operadic bar construction, replaces trees by stable graphs:
\begin{equation}\label{eq:preface-feynman-sum}
F\cO(g,n)
\;=\;
\bigoplus_{\Gamma\in\mathsf{Gr}^{\mathrm{st}}_{g,n}}
\frac{\hbar^{g(\Gamma)}}{|\operatorname{Aut}(\Gamma)|}
\;\bigotimes_{v\in V(\Gamma)}\cO\bigl(g(v),\mathrm{In}(v)\bigr).
\end{equation}
The \emph{differential} contracts one edge at a time:
\[
d_{F\cO}(\gamma_\Gamma)
\;=\;
\sum_{e \in E(\Gamma)}
(-1)^{\epsilon(e)}\,
\xi_e(\gamma_{\Gamma/e}),
\]
where $\gamma_{\Gamma/e}$ denotes the image of $\gamma_\Gamma$
under edge contraction along~$e$, applying the gluing map
$\xi_e$ at edge~$e$.  The sign $(-1)^{\epsilon(e)}$ is determined by
the position of~$e$ in the orientation
line~$\det(E(\Gamma))$.  The identity $d^2 = 0$ is the
codimension-$2$ boundary cancellation: each pair of edges can be
contracted in either order, and the two terms cancel by the
associativity of~$\cO$.  At genus~$0$ (trees, $\hbar^0=1$),
$F\cO$ reduces to the operadic bar construction.

When $\cO=\Com$ (the \emph{commutative modular operad}, with
$\Com(g,n)=\Bbbk$ for every stable $(g,n)$), the Feynman transform
becomes the \emph{stable graph complex}
\begin{equation}\label{eq:preface-fcom}
F\Com(g,n)
\;=\;
\bigoplus_{\Gamma\in\mathsf{Gr}^{\mathrm{st}}_{g,n}}
\frac{\hbar^{g(\Gamma)}}{|\operatorname{Aut}(\Gamma)|}
\;\det\bigl(E(\Gamma)\bigr)\!:
\end{equation}
a chain complex spanned by stable graphs, graded by edge number,
with edge-contraction differential.

Replacing $\Com$ by $\Ass$ equips each vertex with a cyclic
ordering of its half-edges, and $F\!\Ass$ sums over
\emph{ribbon graphs} (fatgraphs) whose genus is the genus of
the thickened embedding surface:
\[
\mathfrak{g}^{\mathrm{mod},\Eone}_\cA
\;:=\;
\Hom_\Sigma\bigl(F\!\Ass,\;
\End_{\barB^{\Eone}(\cA)}\bigr).
\]
This is the algebraic form of the 't~Hooft $1/N$ expansion.
Setting $\hbar = 1/N^2$, the MC element
$\Theta^{\Eone}_\cA \in \operatorname{MC}
(\mathfrak{g}^{\mathrm{mod},\Eone}_\cA)$
is the genus expansion
$\sum_g N^{2-2g}F_g(\lambda)$ summed over ribbon graphs weighted
by the genus of their embedding surface.
Planar diagrams ($g=0$) dominate at large~$N$; the shadow
obstruction tower $\Theta^{E_1,\le r}$ is the truncation to
$r$-point 't~Hooft interactions.
The coinvariant map $F\!\Ass \to F\Com$ projects
$\Theta^{E_1} \mapsto \Theta$:
colour-ordered amplitudes to full amplitudes.

\subsection*{3.3.\enspace The convolution ladder}

The Feynman transform is a skeleton; a chiral algebra provides
the flesh.  The pairing of cooperadic skeleton with operadic
algebra is a convolution dg Lie algebra whose Maurer--Cartan
elements classify algebraic structures:
\begin{equation}\label{eq:preface-convolution-ladder}
\begin{array}{r@{\;:\quad}c@{\qquad}l}
\text{algebra} &
  \Hom\bigl(B(\Ass),\,\End_A\bigr)
  & A_\infty\text{-structures on }A \\[4pt]
\text{operad} &
  \Hom_\Sigma\bigl(B(\cP),\,\End_V\bigr)
  & \cP_\infty\text{-structures on }V \\[4pt]
\text{modular} &
  \Hom_\Sigma\bigl(F\cO,\,\End_{\barB(\cA)}\bigr)
  & \text{consistent genus towers}
\end{array}
\end{equation}
The first line is the Hochschild cochain complex; the
second is the deformation complex of Ginzburg--Kapranov;
the third, with $\cO=\Com$, governs this monograph.  The
\emph{Lie bracket} in each case is the convolution bracket:
for $f, g$ in the convolution hom,
\[
[f, g]
\;:=\;
\mu \circ (f \otimes g) \circ \Delta
\;-\;
(-1)^{|f||g|}\,
\mu \circ (g \otimes f) \circ \Delta,
\]
where $\Delta$ is the cooperadic decomposition and $\mu$ is
operadic composition.

Given a chiral algebra~$\cA$ with bar complex~$\barB(\cA)$,
the \emph{modular convolution dg Lie algebra}
\begin{equation}\label{eq:preface-gmod}
\gAmod
\;:=\;
\Hom_\Sigma\bigl(F\Com,\;\End_{\barB(\cA)}\bigr)
\end{equation}
is the $\cO=\Com$ instance of the modular line.  The bar complex
is an $F\Com$-algebra; the universal MC element $\Theta_\cA :=
D_\cA - d_0$ is bar-intrinsic ($D_\cA^2 = 0$ implies
$\Theta_\cA \in \operatorname{MC}$), and the five main
theorems are projections of~$\Theta_\cA$ (\S4).

\subsection*{3.4.\enspace Homotopy chiral algebras}

The chiral bracket is defined on affine opens; on overlaps, the
Borcherds identity holds only up to higher homotopies.
Fix an affine cover $\mathfrak U=\{U_\alpha\}_{\alpha\in I}$ of the
curve~$X$.  The \v{C}ech totalisation
\[
A^{\mathrm{ch}}_\infty
\;:=\;
\check C^\bullet(\mathfrak U;\cA)
\;=\;
\bigoplus_{p\ge 0}
\prod_{\alpha_0<\cdots<\alpha_p}
\cA(U_{\alpha_0}\cap\cdots\cap U_{\alpha_p})
\]
carries a $\mathrm{Ch}_\infty$-algebra structure
(Malikov--Schechtman): the chiral bracket, restricted to
overlaps of affine opens, satisfies the Jacobi identity only up to
a coherent system of higher homotopies, the \emph{secondary
Borcherds operations}
\[
F_n\colon\mathrm{Lie}(n)\otimes\cY(n)
\;\longrightarrow\;
P_n(\{A^{\mathrm{ch}}_\infty\},\,A^{\mathrm{ch}}_\infty),
\qquad n\ge 2,
\]
encoding derived OPE data at all pole orders.  The operation~$F_2$
is the chiral bracket itself.

The Jacobiator homotopy~$F_3$ arises from the moduli space
$\overline{\cM}_{0,4}\cong\mathbb P^1$: its boundary
$\partial\overline{\cM}_{0,4}$ consists of three points,
corresponding to three OPE compositions.  The chain
$C_1(\overline{\cM}_{0,4})\cong\Bbbk$ provides a
$1$-chain whose boundary is the alternating sum of the three
vertices; $F_3$ is the chain homotopy witnessing their equality
up to~$d$.

Higher operations $F_n$ are determined inductively by
$C_*(\overline{\cM}_{0,n+1})$ through the homotopy transfer
theorem applied to the inclusion
$H_*(\overline{\cM}_{0,n+1})
\hookrightarrow
C_*(\overline{\cM}_{0,n+1})$.
The cellular chain complex of $\overline{\cM}_{0,n+1}$ is
the chain complex of the associahedron~$K_n$; the transferred
operations $F_n$ are the $A_\infty$-corrections at arity~$n$.

Different affine covers yield quasi-isomorphic
$\mathrm{Ch}_\infty$-algebras: the \v{C}ech-to-derived spectral
sequence degenerates for $\mathcal D$-modules on smooth curves.
The strict chiral algebra~$\cA$ is the $H^0$ of this structure: the
primitive local object is the homotopy chiral algebra; the strict
algebra appears after rectification.  The rectification theorem
(Vallette): the homotopy category of $\mathrm{Ch}_\infty$-algebras
is equivalent to the homotopy category of strict chiral algebras,
via the bar-cobar adjunction at the operadic level.

\subsection*{3.5.\enspace Logarithmic Fulton--MacPherson spaces}

The identity $D^2=0$ at higher genus requires the boundary of
the compactified configuration space to be normal-crossings:
every codimension-two stratum must be a transverse intersection
of exactly two codimension-one strata.  The classical FM
compactification does not achieve this in the presence of
a divisor; the logarithmic extension (Mok) does.

For a simple-normal-crossing pair $(X,D)$ with $D$ a reduced
divisor, the \emph{logarithmic Fulton--MacPherson
compactification} $\mathrm{FM}_n(X\mathbin{|}D)$ is a smooth
variety with corners, compactifying the open configuration space
$\mathrm{Conf}_n(X\setminus D)$, whose boundary stratification is
indexed by \emph{rigid planted-forest types}.

A planted-forest type~$\rho$ consists of: a partition of
$\{1,\dots,n\}$ into blocks; for each block, a rooted tree recording
the collision hierarchy; a grid-depth assignment
$w_v\in\mathbb Z_{\ge 0}$ to each vertex, recording the order of
tangency with~$D$.  The ordinary FM compactification is the case
$D=\emptyset$: all grid depths are zero.

The tropicalisation:
\[
\mathrm{Trop}\bigl(\mathrm{FM}_n(\mathbb C\mathbin{|}D)\bigr)
\;=\;
G_{\mathrm{pf}},
\]
the planted-forest poset with grafting as partial order.  Grafting
of planted forests corresponds to degeneration of the underlying
curve; the poset structure is the face poset of the boundary
stratification.

The chain complexes
$\cC^{\log\mathrm{FM}}_{X|D}(n)
:=C_\bullet(\mathrm{FM}_n(X\mathbin{|}D))$ carry cooperadic
structure: the inclusion of boundary strata defines cocomposition
maps.  The \emph{graphwise cocomposition} for a rigid stable
graph~$\Gamma$:
\[
\Delta^{\log}_\Gamma
\;:=\;
(\nu_\Gamma)_*\circ\mathrm{Res}_{D^{\log}_\Gamma}
\;\colon\;
\cC^{\log\mathrm{FM}}_{\mathrm{mod}}(g,n)
\;\longrightarrow\;
\bigotimes_{v\in V(\Gamma)}
\cC^{\log\mathrm{FM}}_{\mathrm{mod}}\bigl(g(v),\,
I_v\sqcup H(v)\bigr)[-|E_{\mathrm{int}}(\Gamma)|],
\]
where $\nu_\Gamma$ is the Mok birational correspondence on the
normalised boundary stratum and
$\mathrm{Res}_{D^{\log}_\Gamma}$ is the residue along the
logarithmic boundary divisor.

\subsection*{3.6.\enspace The convolution $L_\infty$-algebra}

A twisting morphism $\alpha\colon\cC\to\cP$ between a cooperad
and an operad satisfies
$\partial\alpha+\alpha\star\alpha=0$; it twists the
composition product $\cC\circ_\alpha\cP$ into a differential that
squares to zero.  The universal deformation machine built from
$\alpha$ is the convolution
$sL_\infty$-algebra $\mathrm{hom}_\alpha(\cC,\cP)$
(Robert-Nicoud--Wierstra, after Vallette):
\[
\mathrm{hom}_\alpha(\cC,\cP)
\;=\;
\prod_{n\ge 1}
\mathrm{Hom}_{\Sigma_n}\bigl(\cC(n),\,\cP(n)\bigr).
\]
The $L_\infty$ brackets:
\[
\ell_k(f_1,\dots,f_k)(x)
\;=\;
\sum
\gamma_{\cP}\!\bigl(
(\alpha\circ 1)(f_1\otimes\cdots\otimes f_k)\,
\Delta^{(k)}_{\cC}(x)
\bigr),
\]
where $\Delta^{(k)}_{\cC}$ is the iterated cooperadic decomposition
and $\gamma_{\cP}$ is the operadic composition.  The bracket $\ell_1$
is the convolution differential; $\ell_2$ is the Lie bracket of
coderivations; higher brackets $\ell_k$, $k\ge 3$, arise from the
higher cooperadic decompositions and are nontrivial whenever $\cC$
is not concentrated in a single arity.

The $sL_\infty$-algebra $\mathrm{hom}_\alpha(\cC,\cP)$ extends to
$\infty$-morphisms in \emph{either slot separately}: a
quasi-isomorphism $\cC\xrightarrow{\sim}\cC'$ of cooperads induces
an $\infty$-quasi-isomorphism of convolution algebras.  But
\emph{not both simultaneously}: the convolution is functorial in
each slot, but does not assemble into a bifunctor on the
$\infty$-categorical level.  This is a structural constraint, not
a technical gap: the MC3 categorical lift of bar-cobar duality
must proceed one slot at a time.

The bar-cobar adjunction $B_\kappa\dashv\Omega_\kappa$ is a
\emph{Quillen equivalence} (Vallette): different
$\mathrm{Ch}_\infty$ presentations yield quasi-isomorphic
deformation theories.

\subsection*{3.7.\enspace Assembly}

The \emph{modular convolution $L_\infty$-algebra}:
\begin{equation}\label{eq:pref-convolution}
\mathfrak g^{\mathrm{mod}}_\cA
\;:=\;
\Convinf\!\Bigl(
\cC^{\log\mathrm{FM}}_{\mathrm{mod}},\;
\mathrm{End}_{\mathrm{Ch}_\infty}(A^{\mathrm{ch}}_\infty)
\Bigr).
\end{equation}
The cooperadic input is the log-FM chain cooperad
$\cC^{\log\mathrm{FM}}_{\mathrm{mod}}$: the modular cooperad of
chain complexes on logarithmic FM spaces, with cocomposition given
by the graphwise maps~$\Delta^{\log}_\Gamma$.  The operadic input
is the $\mathrm{Ch}_\infty$-endomorphism operad of the homotopy
chiral algebra~$A^{\mathrm{ch}}_\infty$.

The \emph{strict chart}: the coderivation dg Lie algebra
\[
\mathfrak g^{\mathrm{mod},\log}_\cA(\mathfrak U)
\;:=\;
\Convstr\!\bigl(
\cC^{\log\mathrm{FM}}_{\mathrm{mod}},\;
\mathrm{End}_{\mathrm{Ch}_\infty}(A^{\mathrm{ch}}_\infty)
\bigr)
\]
is the \emph{strict model}: classical Lie bracket, no higher
$L_\infty$ brackets, but the \emph{same} Maurer--Cartan moduli.
MC elements in the strict chart are in natural bijection with MC
elements in the full convolution $L_\infty$-algebra.  The strict
chart suffices for constructing the universal MC element and
computing shadows; the full $L_\infty$ structure is needed for
transfer, formality, and gauge equivalence.

\subsection*{3.8.\enspace The six-component differential}

The total differential on the modular bar complex has six
components, whose mutual cancellations constitute $D^2=0$:
\[
D^{\log}_{\mathrm{mod}}
\;=\;
\underbrace{d_{\check C}}_{\text{\v Cech}}
\;+\;
\underbrace{d_{\mathrm{Ch}_\infty}}_{\text{transferred chiral}}
\;+\;
\underbrace{d_{\mathrm{coll}}}_{\text{collision}}
\;+\;
\underbrace{d_{\mathrm{sew}}}_{\text{separating}}
\;+\;
\underbrace{d_{\mathrm{pf}}}_{\text{planted-forest}}
\;+\;
\underbrace{\hbar\,\Delta}_{\text{genus-raising}}.
\]
The collision differential $d_{\mathrm{coll}}$ contracts an internal
edge of a stable graph~$\Gamma$ by extracting the OPE residue at
the corresponding codimension-one collision stratum.  The separating
differential $d_{\mathrm{sew}}$ implements the clutching morphism,
gluing two lower-genus curves along a node.  The genus-raising
operator $\hbar\Delta$ self-glues a curve at two marked points,
increasing $b_1$ by one.  The planted-forest differential
$d_{\mathrm{pf}}$ corrects for the genuinely logarithmic boundary
strata, entering as a push--pull correspondence:
\[
d_{\mathrm{pf}}
\;=\;
\sum_{\rho\in\mathsf{PF}^{\mathrm{rig}}}
\epsilon_\rho\,
(\kappa_\rho)_*\,\mathrm{pr}_\rho^*\otimes\mu_\rho\,,
\]
where $\mathrm{pr}_\rho^*$ pulls back chains to the planted-forest
stratum and $(\kappa_\rho)_*$ pushes forward along the Mok
birational correspondence.

\subsection*{3.9.\enspace $D^2=0$}

The nilpotence $(D^{\log}_{\mathrm{mod}})^2=0$ holds by three
cancellation mechanisms:
\begin{enumerate}[label=(\roman*)]
\item \emph{Vertex labels}: $d_{\check C}^{\,2}=0$,
$d_{\mathrm{Ch}_\infty}^{\,2}=0$, and their anticommutator
$\{d_{\check C},d_{\mathrm{Ch}_\infty}\}=0$ by the
$\mathrm{Ch}_\infty$-algebra axioms.
\item \emph{One-edge expansions}: the anticommutator of a
vertex differential with an edge-contraction differential
vanishes by the Leibniz rule for coderivations.
\item \emph{Codimension-two cancellation}: the sum of all
products of two distinct edge-contraction operators vanishes.
The codimension-two boundary of
Mok's log-FM compactification is a normal-crossings divisor
(Mok, Theorem~3.3.1), and each codimension-two face appears as
the intersection of exactly two codimension-one faces.  The
two boundary contributions cancel with opposite orientations:
$\partial^2=0$ for the boundary of a manifold with corners,
lifted through the residue correspondence.
\end{enumerate}
The third mechanism is the nontrivial one.  The normal-crossings
property of Mok's compactification is essential: it guarantees
that every codimension-two intersection is transverse and
contributes to exactly one cancelling pair.

\begin{remark}[Homotopy types in three stages]
\label{rem:pref-homotopy-types}
The three levels of the convolution ladder mirror three levels of
homotopy theory.  \emph{Over a point}: the bar construction
$B(A) = (T^c(s^{-1}\bar A), d_B)$ determines the homotopy type of
an associative algebra; formality ($m_k=0$ for $k \ge 3$)
corresponds to intrinsic formality; the Massey products are the
obstructions.  \emph{Over a manifold}: the Sullivan minimal model
$(\Lambda V, d)$ determines the rational homotopy type; K\"ahler
manifolds are formal (Deligne--Griffiths--Morgan--Sullivan);
the Heisenberg algebra is the vertex-algebraic analogue, with
shadow obstruction tower terminating at arity~$2$.
\emph{Over a curve at all genera}: the $A_\infty$-products
$m_k$ become modular $L_\infty$-brackets
$\ell_k^{(g)}$; formality becomes the Gaussian archetype
($\ell_k^{(g)}=0$ for $k \ge 3$ at all~$g$); the shadow
obstruction tower is the modular Massey product tower, classifying
the complexity of the modular homotopy type.
The four classes G, L, C, M (Section~6) correspond to four levels
of modular formality, determined by the vanishing pattern of the
higher brackets~$\ell_k^{(g)}$.
\end{remark}


% ====================================================================
\section*{4.\quad The universal Maurer--Cartan element}
% ====================================================================

\medskip

\noindent
The modular convolution algebra $\gAmod$ has a distinguished
Maurer--Cartan element $\Theta_\cA$ that requires no equation to
solve: it is extracted from the total bar differential, which
already squares to zero.  The decomposition $D_\cA = d_0 + \Theta_\cA$
splits the differential into its tree-level local part and the
remainder; the MC equation for $\Theta_\cA$ is a formal consequence
of $D_\cA^2 = 0$.  Every invariant in this monograph is a projection
of $\Theta_\cA$.

\subsection*{4.1.\enspace $\Gamma$-amplitudes and Taylor coefficients}

Each stable graph produces a multilinear operation on the convolution
algebra from three ingredients: chiral operations at vertices,
propagators along edges, and geometric factors from the log-FM
strata.

For a stable graph $\Gamma\in\mathsf{Gr}^{\mathrm{st}}$ with
$|V(\Gamma)|=k$ vertices and elements $f_1,\dots,f_k\in\gAmod$
assigned to those vertices, the \emph{$\Gamma$-amplitude}:
let $\mu_\Gamma\colon\bigotimes_{v\in V(\Gamma)}A^{\otimes n(v)}\to A$
denote the graph composition; at each vertex~$v$ of genus~$g(v)$ and
valence~$n(v)$, place the transferred chiral operation
$m^{(h)}_{g(v),n(v)}$; along each internal edge, contract with the
propagator~$P_\cA=\iota\circ\pi$ (homotopy retraction of the
augmentation ideal to its cohomology).  Write
$\alpha^{\mathrm{mod}}_\Gamma$ for the component of the modular
twisting morphism on~$\Gamma$, and
$\Delta^{\log}_\Gamma$ for the cooperadic cocomposition.
The $\Gamma$-amplitude:
\[
\ell_\Gamma(f_1,\dots,f_k)
\;=\;
\mu_\Gamma\circ
(\alpha^{\mathrm{mod}}_\Gamma\otimes f_1\otimes\cdots\otimes f_k)
\circ\Delta^{\log}_\Gamma.
\]
This is a $k$-linear map
$\ell_\Gamma\colon(\gAmod)^{\otimes k}\to\gAmod$ of cohomological
degree $2-2g(\Gamma)-k$.

The modular $L_\infty$ structure on~$\gAmod$ has Taylor coefficients
obtained by summing over all stable graphs of fixed type:
\[
\ell_k^{(g)}
\;=\;
\sum_{\substack{
\Gamma\in\mathsf{Gr}^{\mathrm{st}},\;|V(\Gamma)|=k\\
b_1(\Gamma)+\sum_v g(v)=g}}
\frac{1}{|\operatorname{Aut}(\Gamma)|}\;\ell_\Gamma.
\]
Here $\ell_k^{(g)}$ is a symmetric $k$-linear bracket of genus~$g$.
The full modular $L_\infty$ structure is the formal sum
$\ell=\sum_{g\ge 0}\sum_{k\ge 1}\hbar^g\,\ell_k^{(g)}/k!$, and the
homotopy Jacobi identity is encoded by $D^2=0$.

\medskip

\noindent\textbf{Explicit low-order brackets.}\enspace
At tree level ($g=0$), the binary bracket is the Lie bracket on
cohomology transferred through the homotopy retraction data
$(\pi,\iota,h)$:
\[
\ell_2^{(0)}(\alpha,\beta)
\;=\;
\pi\bigl[m_2\bigl(\iota(\alpha),\iota(\beta)\bigr)\bigr].
\]
The ternary bracket sums over three binary trees (one for each
channel $s$, $t$, $u$):
\begin{align*}
\ell_3^{(0)}(\alpha,\beta,\gamma)
&=\pi\Bigl(
m_2\bigl(h\,m_2(\iota\alpha,\iota\beta),\,\iota\gamma\bigr)
+m_2\bigl(\iota\alpha,\,h\,m_2(\iota\beta,\iota\gamma)\bigr)\\
&\qquad\quad
+m_2\bigl(h\,m_2(\iota\alpha,\iota\gamma),\,\iota\beta\bigr)
\Bigr).
\end{align*}
The homotopy~$h$ propagates the intermediate state through~$P_\cA$
before re-composing: the classical homological perturbation lemma
in its $L_\infty$ incarnation.

The first genus-raised bracket is unary:
\[
\ell_1^{(1)}(\alpha)
\;=\;
\operatorname{tr}(P_\cA\circ\alpha\circ
\delta^{\mathrm{ns}})
\;=\;
\sum_{i}
\langle e^i,\,\alpha(e_i)\rangle_{\mathrm{Sh}},
\]
where $\{e_i\}$, $\{e^i\}$ is a dual basis for the cyclic pairing
on~$H^*(\bar\cA)$ and $\delta^{\mathrm{ns}}$ is the non-separating
node operator.  Geometrically, this is the BV operator: the unique
genus-one graph with one vertex and one self-loop, the trace over the
propagator contracting the two half-edges at the non-separating node
of~$\overline{\cM}_{1,1}$.

\subsection*{4.2.\enspace The bar-intrinsic construction}

The total bar differential
$D_\cA\colon\mathrm{B}^{\mathrm{mod}}(\cA)\to
\mathrm{B}^{\mathrm{mod}}(\cA)$ acts on the modular bar complex.
Decompose $D_\cA=d_0+\Theta_\cA$, where
$d_0=d_{\mathrm{int}}+[\tau,-]$ is the tree-level local part: the
internal differential of~$\cA$ plus the commutator with the
tautological twisting morphism~$\tau$.  The \emph{universal
Maurer--Cartan element} is the remainder:
\begin{equation}\label{eq:pref-theta}
\Theta_\cA
\;:=\;
D_\cA-d_0
\;=\;
\sum_{\Gamma\in\mathsf{Gr}^{\mathrm{st}}}
\frac{W^{\log\mathrm{FM}}_\Gamma}{|\operatorname{Aut}(\Gamma)|}
\otimes\Phi^\cA_\Gamma
\;\in\;
\operatorname{MC}(\gAmod).
\end{equation}
Here $W^{\log\mathrm{FM}}_\Gamma\in
H^*(\overline{\cM}^{\log}_{g(\Gamma),n(\Gamma)})$ is the log-FM
weight (the fundamental class of the boundary stratum indexed
by~$\Gamma$), and
$\Phi^\cA_\Gamma\in\operatorname{Hom}(\bar\cA^{\otimes
n(\Gamma)},\Bbbk)$ is the chiral amplitude on~$\Gamma$.

The MC property is immediate: $D_\cA^{\,2}=0$ by the geometric fact
that $\overline{\cM}_{g,n}^{\log}$ has normal-crossings boundary;
$d_0^2=0$ by construction; expanding $(d_0+\Theta_\cA)^2=0$ gives
\[
d_0\Theta_\cA+\Theta_\cA\,d_0
+\Theta_\cA^2=0
\qquad\Longleftrightarrow\qquad
[d_0,\Theta_\cA]+\tfrac12[\Theta_\cA,\Theta_\cA]=0.
\]
No equation is solved; the MC property is a formal consequence of
$\partial^2=0$ on the moduli space of curves.  No simple-Lie
restriction, no convergence hypothesis, no choice of basis.

On the $E_1$ side, the same decomposition applied to the ordered bar
differential $D^{E_1}_\cA$ on $T^c(s^{-1}\bar\cA)$ produces the
ordered MC element $\Theta^{E_1}_\cA \in
\operatorname{MC}(\mathfrak{g}^{\mathrm{mod},E_1}_\cA)$.  The
coinvariant projection sends $\Theta^{E_1}_\cA \mapsto \Theta_\cA$:
the modular MC element is the $\Sigma_n$-averaged shadow of the
ordered one.  At arity~$2$, this recovers
$\mathrm{av}(r(z)) = \kappa(\cA)$.

\subsection*{4.3.\enspace Universal modular twisting morphism}

Twisting morphisms classify compatible deformations of the
bar complex across all genera simultaneously.  The MC space of
the modular convolution $L_\infty$-algebra is the space of
modular twisting morphisms:
\[
\operatorname{MC}_\bullet\!\bigl(\gAmod\bigr)
\;\simeq\;
\operatorname{Tw}_{\alpha^{\mathrm{mod}}}\!\Bigl(
\cC^{\log\mathrm{FM}}_{\mathrm{mod}},\,
\operatorname{End}_{\mathrm{Ch}_\infty}(A^{\mathrm{ch}}_\infty)
\Bigr).
\]
Under this identification, $\Theta_\cA$ is the \emph{universal
modular twisting morphism}: a twisting
morphism $\alpha\colon\cC\to\cP$ determines a twisted composition
product $\cC\circ_\alpha\cP$ whose differential squares to zero if
and only if~$\alpha$ satisfies the MC equation.  $\Theta_\cA$ is the
unique twisting morphism that simultaneously twists the modular bar
complex and intertwines with the log-FM boundary structure.

The principal-bundle analogy is exact.  The modular bar
bundle $\mathrm{B}^{\log}_{\mathrm{mod}}(\cA)\to
\overline{\cM}_{g,n}$ is the principal bundle; $\Theta_\cA$ is the
connection; the MC equation is the flatness condition; the shadow
tower (Section~6) is the Chern--Weil theory.  The structure group is
the pro-unipotent MC gauge group
$\exp(\mathfrak g^{\mathrm{mod},0}_\cA)$.

\subsection*{4.4.\enspace Tridegree and the depth filtration}

Each coefficient of~$\Theta_\cA$ carries a natural tridegree
\[
(g,n,d)
\;=\;
(\text{loop genus},\;\text{external arity},\;
\text{log boundary depth}),
\]
where $d(\rho)=\#E_{\mathrm{sep}}(\Gamma_\rho)+
\#E_{\mathrm{ns}}(\Gamma_\rho)+\operatorname{depth}_{\mathrm{pf}}
(\Gamma_\rho)$ is the codimension of the log-FM stratum indexed by
the degeneration pattern~$\rho$.

The depth filtration
$F^d\gAmod=\bigoplus_{d'\ge d}(\gAmod)_{*,*,d'}$ makes
the convolution algebra complete:
\[
\gAmod\;=\;\varprojlim_{d}\;\gAmod/F^d\gAmod.
\]
The inverse limit
$\Theta_\cA=\varprojlim_r\Theta_\cA^{\le r}$ converges
unconditionally: no analytic input, no growth estimate, no
restriction on~$\cA$.  The convergence is purely algebraic,
guaranteed by the completeness of the depth filtration.

At fixed $(g,n)$, only finitely many stable graph types~$\Gamma$
contribute (by stability: $2g(\Gamma)-2+n(\Gamma)>0$ forces a finite
graph set).  The depth filtration refines this: at depth~$d$,
the stratum has codimension~$d$ in $\overline{\cM}_{g,n}^{\log}$,
and the amplitude $\Phi^\cA_\Gamma$ has propagator insertions
bounded by~$d$.

\subsection*{4.5.\enspace The weight filtration}

Decompose the MC equation by internal-edge count:
\[
\Theta_\cA\;=\;\sum_{r\ge 1}\Theta^{[r]}_\cA,
\qquad
\Theta^{[r]}_\cA\;=\;
\sum_{\substack{\Gamma\in\mathsf{Gr}^{\mathrm{st}}\\
|E_{\mathrm{int}}(\Gamma)|=r}}
\frac{W^{\log\mathrm{FM}}_\Gamma}{|\operatorname{Aut}(\Gamma)|}
\otimes\Phi^\cA_\Gamma.
\]
The MC equation decomposes by weight into an inductive tower:
\begin{alignat*}{2}
&\text{Weight }1{:}\quad
&&[D_{\mathrm{loc}},\Theta^{[1]}_\cA]=0,\\[4pt]
&\text{Weight }2{:}\quad
&&[D_{\mathrm{loc}},\Theta^{[2]}_\cA]
+\tfrac12[\Theta^{[1]}_\cA,\Theta^{[1]}_\cA]=0,\\[4pt]
&\text{Weight }N{:}\quad
&&[D_{\mathrm{loc}},\Theta^{[N]}_\cA]
+\tfrac12\sum_{\substack{i+j=N\\i,j\ge 1}}
[\Theta^{[i]}_\cA,\Theta^{[j]}_\cA]=0.
\end{alignat*}
At each stage, $[D_{\mathrm{loc}},\Theta^{[N]}]$ is determined by
the quadratic error from previous stages.  The obstruction to
extending from weight~$N-1$ to weight~$N$ lies in
$H^2(\gAmod/F^{N+1},D_{\mathrm{loc}})$.  The bar-intrinsic
construction solves all stages simultaneously.

\medskip

\noindent\textbf{Weight~$1$: boundary divisors.}\enspace
The weight-one component $\Theta^{[1]}_\cA$ lives on boundary
divisors of moduli spaces.  In degree $(0,4)$: the three boundary
points of $\overline{\cM}_{0,4}\cong\mathbb P^1$:
\[
\Theta^{[1]}_\cA\big|_{(0,4)}
=[\delta_s]\otimes\Phi_s
+[\delta_t]\otimes\Phi_t
+[\delta_u]\otimes\Phi_u.
\]
The weight-one MC equation at $(0,4)$ reads
$\Phi_s+\Phi_t+\Phi_u=F_3$,
the Malikov--Schechtman Jacobiator homotopy: the Jacobiator is a
boundary in the bar complex.

In degree $(1,1)$: the single boundary divisor
$\delta_{\mathrm{irr}}\subset\overline{\cM}_{1,1}$ gives
\[
\Theta^{[1]}_\cA\big|_{(1,1)}
=[\delta_{\mathrm{irr}}]\otimes\Delta,
\]
where $\Delta=\sum_i e_i\otimes e^i$ is the Casimir element of the
cyclic pairing.

In degree $(0,5)$: the weight-one component restricts to the ten
boundary divisors of $\overline{\cM}_{0,5}$ (the Petersen graph),
and the MC equation recovers the pentagon identity.  The entire
$A_\infty$ hierarchy is the tree-level, genus-zero projection of a
single MC element.

\subsection*{4.6.\enspace The all-genus recursion}

Project the MC equation to fixed genus~$g$ and arity~$n$:
\[
d\Theta^{(g,n)}_\cA
+\tfrac12
\!\!\!
\sum_{\substack{g_1+g_2=g\\n_1+n_2=n+2}}
\!\!\!
\ell_2^{(0)}(\Theta^{(g_1,n_1)}_\cA,\Theta^{(g_2,n_2)}_\cA)
+\Delta_{\mathrm{ns}}\,\Theta^{(g-1,n+2)}_\cA
+\sum_{k\ge 3}\frac{1}{k!}
\ell_k^{(\mathrm{tree})}(\Theta_\cA,\dots,\Theta_\cA)^{(g,n)}
=0.
\]
The four terms: \emph{internal differential} at fixed $(g,n)$;
\emph{separating clutching} over splittings $g_1+g_2=g$,
$n_1+n_2=n+2$; \emph{non-separating BV operator} lowering arity
by~$2$ and raising genus by~$1$; \emph{planted-forest corrections}
from higher brackets $\ell_k^{(\mathrm{tree})}$ for $k\ge 3$.

This recursion determines $\Theta^{(g,n)}_\cA$ inductively from
lower $(g',n')$ in lexicographic order.  The base cases are
$\Theta^{(0,3)}_\cA$ (the chiral bracket) and
$\Theta^{(0,4)}_\cA$ (the Jacobiator data from weight~$1$).
The recursion terminates at each $(g,n)$: only finitely many stable
graphs contribute.

\medskip
\noindent\textbf{The clutching identity.}\enspace
On each rigid boundary stratum~$\rho$ of the log-FM
compactification, the restriction of~$\Theta_\cA$ factors as a
tensor product over vertices of the dual graph:
\[
\Theta_\cA\big|_\rho
\;=\;
(\kappa_\rho)_*\!\left(
\bigotimes_{v\in V(S_\rho)}\Theta_{\cA,v}
\right).
\]
The amplitude on the degenerate curve equals the product
of amplitudes on the irreducible components, sewn by propagators
along the nodes: the MC-level incarnation of the sewing axiom.

\bigskip
\noindent
Every invariant in this preface is a projection of~$\Theta_\cA$:
the curvature~$\kappa$, the spectral discriminant~$\Delta_\cA$,
the shadow metric~$Q_L$, the sewing amplitudes, the $\hat A$-genus,
the Fredholm determinant, and the quartic contact
invariant~$Q^{\mathrm{contact}}$.
The five main theorems are five projections of~$\Theta_\cA$.

\medskip
\begin{center}
\small
\renewcommand{\arraystretch}{1.15}
\begin{tabular}{lll}
\textbf{Projection of $\Theta_\cA$}
  & \textbf{Invariant} & \textbf{Section} \\
\hline
Scalar ($r=2$, genus tower)
  & $\kappa(\cA)$, $F_g$, $\hat A$-genus
  & \S2 \\[2pt]
Binary genus-$0$ residue
  & classical $r$-matrix, Yangian seed
  & \S7 \\[2pt]
Binary Verdier dual
  & complementarity $Q_g \oplus Q_g^!$
  & \S2 \\[2pt]
Bar-cobar counit
  & inversion $\Omega(\barB(\cA))\xrightarrow{\sim}\cA$
  & \S1 \\[2pt]
Chiral Hochschild
  & $\mathrm{ChirHoch}^*(\cA)$, deformation ring
  & \S9 \\[2pt]
Spectral discriminant
  & $\Delta_\cA(x)$, genus-$1$ Hessian
  & \S5 \\[2pt]
Cubic shadow ($r=3$)
  & gauge-trivial obstruction
  & \S6 \\[2pt]
Quartic resonance ($r=4$)
  & $\mathfrak Q^{\mathrm{contact}}$, clutching law
  & \S6 \\[2pt]
MC tautological descent
  & shadow CohFT, $R^*(\overline{\cM}_{g,n})$
  & \S5 \\[2pt]
Swiss-cheese colour decomposition
  & $\alpha_T$, six projections (Vol.~II)
  & \S10
\end{tabular}
\end{center}
\medskip

\noindent
The existence of~$\Theta_\cA$ is encoded by the bar-cobar
adjunction: Theorem~A constructs the arena in which~$\Theta_\cA$
lives.  Bar-cobar inversion proves that~$\Theta_\cA$ is a complete
invariant: Theorem~B recovers~$\cA$ from~$\Theta_\cA$.  Verdier
duality decomposes~$\Theta_\cA$ into complementary halves: Theorem~C
splits the genus tower into Lagrangian summands.  The leading
coefficient of~$\Theta_\cA$ is the modular characteristic: Theorem~D
extracts the scalar $\kappa(\cA)$ with $\hat A$-genus generating
function.  The coefficient ring of~$\Theta_\cA$ is the chiral
Hochschild complex: Theorem~H identifies the deformation ring.

The finite-order projections form an obstruction tower
$\Theta_\cA^{\leq 2} \to
\Theta_\cA^{\leq 3} \to
\Theta_\cA^{\leq 4} \to \cdots$
whose generating function is algebraic of degree~$2$: on each
primary line~$L$, the weighted shadow generating function
$H(t) = t^2\sqrt{Q_L(t)}$ for a quadratic polynomial~$Q_L$
determined by three seed invariants $(\kappa, \alpha, S_4)$.
The inverse limit
$\Theta_\cA = \varprojlim_r \Theta_\cA^{\leq r}$
is proved to exist and to be automatically Maurer--Cartan
because $D_\cA^2 = 0$.

\medskip
\noindent\textbf{Corollary} (MC replaces the sum over topologies).
On the Koszul locus, the tree-level complex carries a contracting
homotopy (from the PBW degeneration), so the genus-$g$ component
$\Theta^{(g)}$ is determined by induction from $\Theta^{(0)}$.  No
free parameters appear at any genus; distinct choices of strong
deformation retract yield gauge-equivalent MC elements.  The genus
expansion is computed algebraically from the genus-$0$ chiral bracket
by the clutching recursion, not approximated perturbatively.  The
Maloney--Witten problem (negative Fourier coefficients in the
Poincar\'e series) does not arise: the MC equation determines each
$F_g$ uniquely by algebraic induction, without summing over bulk
topologies at fixed genus.
